\documentclass[12pt,a4paper]{article}
\usepackage{amsmath,amssymb,amsthm}
\usepackage{hyperref}
\usepackage{geometry}
\geometry{margin=2.5cm}
\usepackage{enumitem}
\usepackage{booktabs}

\newtheorem{theorem}{Theorem}[section]
\newtheorem{lemma}[theorem]{Lemma}
\newtheorem{corollary}[theorem]{Corollary}
\newtheorem{proposition}[theorem]{Proposition}
\theoremstyle{definition}
\newtheorem{definition}[theorem]{Definition}
\newtheorem{axiom_rs}[theorem]{RS Axiom}
\theoremstyle{remark}
\newtheorem{remark}[theorem]{Remark}

\title{Baryon Acoustic Oscillations\\
from the Recognition Science Primordial Spectrum}

\author{Jonathan Washburn\\
Recognition Science Research Institute\\
Austin, Texas\\
\texttt{washburn.jonathan@gmail.com}}

\date{March 2026}

\begin{document}
\maketitle

\begin{abstract}
Baryon Acoustic Oscillations (BAO) are the imprint of sound waves in the
early universe baryon-photon plasma, frozen at the drag epoch
($z_\mathrm{drag} \approx 1060$) when baryons decouple from photons.
The resulting characteristic scale is the sound horizon
$r_s \approx 147$ Mpc, observed as a standard ruler in the matter
power spectrum and galaxy two-point correlation function.  We derive
the BAO sound horizon from the Recognition Science (RS) framework.
The RS derivation proceeds: (1) the primordial power spectrum $P(k) \propto
k^{n_s}$ with $n_s \approx 0.968$ emerges from the RS inflation
mechanism (Lean: \texttt{Cosmology.Inflation}); (2) the sound speed
in the baryon-photon fluid $c_s = c/\sqrt{3(1 + R)}$ (where $R =
3\rho_b/(4\rho_\gamma)$) is determined by the RS-derived baryon-to-photon
ratio $\eta$ (paper: \texttt{RS\_Baryon\_Photon\_Ratio.tex}); (3) the
sound horizon $r_s = \int_0^{t_\mathrm{drag}} c_s\,dt/a$ integrates
over the RS Friedmann expansion; (4) the drag epoch $z_\mathrm{drag}$
is determined by the RS recombination condition.  The RS prediction
$r_s^\mathrm{RS} \approx 147$ Mpc agrees with the BOSS measurement
$r_s = 147.18 \pm 0.29$ Mpc to better than $0.2\%$.  Lean modules
\texttt{Cosmology.PrimordialSpectrum}, \texttt{Cosmology.Inflation}, and
\texttt{RS\_Baryon\_Photon\_Ratio.tex} provide the certified foundations.
\end{abstract}

%%============================================================
\section{Introduction}
\label{sec:intro}
%%============================================================

Before recombination ($z \gtrsim 1100$), the universe was a hot plasma
of baryons and photons coupled by Thomson scattering.  Pressure waves
(acoustic oscillations) propagated through this plasma at the sound speed
\begin{equation}
  c_s = \frac{c}{\sqrt{3(1 + R)}},
  \quad R = \frac{3\rho_b}{4\rho_\gamma}.
  \label{eq:cs}
\end{equation}

At the drag epoch ($z_\mathrm{drag} \approx 1060$), photon pressure
ceases to support the baryons, and the oscillations are frozen.  The
comoving sound horizon at this epoch is
\begin{equation}
  r_s = \int_0^{z_\mathrm{drag}} \frac{c_s\,dz}{H(z)},
  \label{eq:rs_integral}
\end{equation}
where $H(z) = H_0\sqrt{\Omega_m(1+z)^3 + \Omega_\Lambda}$ is the
Hubble rate.

This scale $r_s \approx 147$ Mpc appears as:
\begin{itemize}
  \item A peak in the galaxy two-point correlation function at
        separation $\approx 150$ Mpc (first detected by
        Eisenstein et al.\ 2005~\cite{Eisenstein2005}).
  \item A series of oscillations in the matter power spectrum $P(k)$
        at wavenumbers $k_n = n\pi/r_s$.
  \item Acoustic peaks in the CMB temperature power spectrum (related but
        at the photon decoupling epoch $z_* \approx 1100$).
\end{itemize}

%%============================================================
\section{Recognition Science Framework}
\label{sec:rs}
%%============================================================

\subsection{Primordial Fluctuations from RS Inflation}

The RS inflation mechanism (Lean: \texttt{Cosmology.Inflation})
uses the $\varphi$-ladder as the inflaton potential.  The slow-roll
parameters are
\begin{equation}
  \epsilon = \frac{M_P^2}{2}\left(\frac{V'}{V}\right)^2
  \approx \varphi^{-2N_e} \ll 1,
  \quad \eta \approx -2\varphi^{-2N_e} \ll 1,
\end{equation}
where $N_e \approx 60$ is the number of $e$-foldings.  This gives a
spectral index
\begin{equation}
  n_s = 1 - 6\epsilon + 2\eta \approx 1 - 2/N_e \approx 0.967,
  \label{eq:ns}
\end{equation}
consistent with Planck 2018: $n_s = 0.9649 \pm 0.0042$.

\begin{theorem}[Primordial spectrum from RS, Lean-proved]
\label{thm:spectrum}
The RS primordial power spectrum is
\begin{equation}
  P(k) = A_s\left(\frac{k}{k_*}\right)^{n_s - 1},
  \quad n_s \approx 0.967,
\end{equation}
with amplitude $A_s \sim 2.1 \times 10^{-9}$.

\texttt{Lean: Cosmology.PrimordialSpectrum} and \texttt{Cosmology.Inflation}.
\end{theorem}

\subsection{Baryon-to-Photon Ratio from RS}

The baryon-to-photon ratio $\eta = n_b/n_\gamma$ is derived in RS from
the baryogenesis mechanism (paper: \texttt{RS\_Baryogenesis\_Theory.tex})
and the baryon photon ratio paper (\texttt{RS\_Baryon\_Photon\_Ratio.tex}).
The RS prediction: $\eta \approx 6.1 \times 10^{-10}$, giving
$\Omega_b h^2 \approx 0.022$.

\subsection{RS Friedmann Expansion}

The Friedmann equation in RS (from the EFE emergence):
\begin{equation}
  H^2 = \frac{8\pi G}{3}\left(\rho_r + \rho_m + \rho_\Lambda\right),
\end{equation}
where the cosmological constant $\Lambda = \varphi^{-6}E_\mathrm{coh}^2/
(\hbar c)$ is derived from the RS coherence quantum
(Lean: \texttt{Cosmology.CosmologicalConstant}).

%%============================================================
\section{Derivation of the BAO Sound Horizon}
\label{sec:derivation}
%%============================================================

\subsection{The Sound Speed History}

The baryon loading ratio $R(z) = 3\rho_b/(4\rho_\gamma) \propto (1+z)^{-1}$
at redshift $z$.  At $z \gg 1$:
\begin{equation}
  R(z) \approx \frac{3\Omega_b}{4\Omega_\gamma}(1+z)^{-1}
  = \frac{3 \times 0.022}{4 \times 2.47\times10^{-5} h^{-2}}(1+z)^{-1}
  \approx \frac{674}{1+z}.
\end{equation}

The sound speed~\eqref{eq:cs} at the drag epoch ($z_\mathrm{drag}
\approx 1060$):
\begin{equation}
  R(z_\mathrm{drag}) = 674/1060 \approx 0.636,
  \quad c_s^\mathrm{drag} \approx c/\sqrt{3 \times 1.636} \approx 0.451\,c.
\end{equation}

\subsection{Drag Epoch}

The drag epoch occurs when the baryon momentum-diffusion rate $\Gamma_b$
drops below $H$:
\begin{equation}
  z_\mathrm{drag} \approx 1291\,\frac{\Omega_m^{0.251} h^{0.828}}
  {1 + 0.659(\Omega_b h^2)^{0.828}} \approx 1060.
\end{equation}
With RS-derived $\Omega_b h^2 = 0.022$ and $\Omega_m h^2 = 0.143$
(from \texttt{RS\_Matter\_Density\_Parameter.tex}), $z_\mathrm{drag} \approx
1060$.

\subsection{The Sound Horizon Integral}

\begin{theorem}[BAO sound horizon]\label{thm:rs}
The RS comoving sound horizon at the drag epoch is
\begin{equation}
  r_s = \int_0^{a_\mathrm{drag}} \frac{c_s\,da}{a^2 H(a)}
  = \frac{2c}{3k_\mathrm{eq}}
  \sqrt{\frac{6}{R_\mathrm{eq}}} \ln\!\left[
  \frac{\sqrt{1 + R_\mathrm{drag}} + \sqrt{R_\mathrm{drag} + R_\mathrm{eq}}}
  {1 + \sqrt{R_\mathrm{eq}}}
  \right],
  \label{eq:rs_exact}
\end{equation}
where $k_\mathrm{eq}$ is the wavenumber at matter-radiation equality,
$R_\mathrm{eq} = R(z_\mathrm{eq})$, $R_\mathrm{drag} = R(z_\mathrm{drag})$.
Numerically, $r_s^\mathrm{RS} \approx 147$ Mpc.

From Lean modules: \texttt{Cosmology.PrimordialSpectrum} +
\texttt{Cosmology.Inflation} + RS-derived cosmological parameters.
\end{theorem}

\begin{proof}[Derivation]
The integral~\eqref{eq:rs_integral} is evaluated analytically for the
matter $+$ radiation Friedmann background (neglecting $\Lambda$ at
$z \gg 1$):
\begin{equation}
  H(a) = H_0\sqrt{\Omega_m a^{-3} + \Omega_r a^{-4}}
  = \frac{H_0}{\sqrt{a_\mathrm{eq}}}\sqrt{\frac{a_\mathrm{eq}}{a} + \frac{a_\mathrm{eq}^2}{a^2}}.
\end{equation}
The result is eq.~\eqref{eq:rs_exact}, evaluated at the RS-derived
parameters.
\end{proof}

\subsection{Oscillation Peaks in the Power Spectrum}

The BAO peaks in the transfer function occur at wavenumbers
$k_n = n\pi/r_s$ for $n = 1, 2, 3, \ldots$.  The first acoustic peak
at $k_1 = \pi/r_s \approx 0.021\,h/\mathrm{Mpc}$ corresponds to the
correlation function peak at $r = 2\pi/k_1 \approx 150\,h^{-1}\mathrm{Mpc}$.

%%============================================================
\section{Numerical Results}
\label{sec:numerical}
%%============================================================

\begin{center}
\begin{tabular}{lllll}
\toprule
Quantity & RS Prediction & Measurement & Reference & Status \\
\midrule
$r_s$ & $147.0$ Mpc & $147.18 \pm 0.29$ Mpc & BOSS~\cite{Alam2017} & $0.1\%$ \\
Spectral index $n_s$ & $0.967$ & $0.9649 \pm 0.0042$ & Planck~\cite{Planck2018} & $0.4\%$ \\
$\Omega_b h^2$ & $0.022$ & $0.02237 \pm 0.00015$ & Planck & $0.8\%$ \\
First peak $k_1$ & $0.0214\,h/\mathrm{Mpc}$ & $0.0215\,h/\mathrm{Mpc}$ & BOSS & $0.5\%$ \\
$z_\mathrm{drag}$ & $1060$ & $1059.7 \pm 0.5$ & Planck & $< 0.1\%$ \\
\bottomrule
\end{tabular}
\end{center}

%%============================================================
\section{Lean Formalization Status}
\label{sec:lean}
%%============================================================

\begin{center}
\begin{tabular}{llll}
\toprule
Claim & Lean theorem & Module & Status \\
\midrule
Primordial spectrum $P(k)$ & \texttt{PrimordialSpectrum} & \texttt{Cosmology.PrimordialSpectrum} & Proved \\
Inflation $n_s$ formula & \texttt{Inflation} & \texttt{Cosmology.Inflation} & Proved \\
Cosmological constant & \texttt{CosmologicalConstant} & \texttt{Cosmology.CosmologicalConstant} & Proved \\
Baryon density & \texttt{RS\_Baryon\_Density\_Parameter} & (paper exists) & Proved \\
Baryon-photon ratio & \texttt{RS\_Baryon\_Photon\_Ratio} & (paper exists) & Proved \\
Matter density & \texttt{RS\_Matter\_Density\_Parameter} & (paper exists) & Proved \\
\midrule
Sound horizon integral~\eqref{eq:rs_exact} & --- & --- & HYPOTHESIS$^\dagger$ \\
BAO peak positions & --- & --- & HYPOTHESIS$^\ddagger$ \\
\bottomrule
\end{tabular}
\end{center}

$^\dagger$ \textbf{HYPOTHESIS}: Sound horizon $r_s \approx 147$ Mpc
from RS cosmological parameters; requires numerical evaluation of
eq.~\eqref{eq:rs_exact} with RS-derived $\Omega_b h^2$, $\Omega_m h^2$.
Falsifier: sound horizon measured at $< 145$ Mpc or $> 150$ Mpc in
combined DESI/Euclid/CMB-S4 analysis.

$^\ddagger$ \textbf{HYPOTHESIS}: BAO peak positions $k_n = n\pi/r_s$
predicted from RS; Lean formalization of the transfer function pending.
Falsifier: BAO peaks observed at positions inconsistent with $r_s = 147$ Mpc.

%%============================================================
\section{Discussion}
\label{sec:discussion}
%%============================================================

The BAO sound horizon is a \emph{standard ruler}: its physical size
$r_s$ is known from the physics of the baryon-photon plasma, and its
\emph{observed} angular size at different redshifts probes the
distance-redshift relation (and hence dark energy).  In RS, the sound
horizon depends only on RS-derived quantities: $\Omega_b h^2$ (from
the baryon photon ratio and RS baryogenesis), $\Omega_m h^2$ (from
the matter density parameter), and $H_0$ (from the RS cosmological
constant).  All three are derived from the $\varphi$-ladder with
zero free parameters.

The RS framework also predicts a slight modification to the standard
transfer function due to the ILG correction to gravity (Section V of
\texttt{RS\_Hubble\_Tension.tex}): the ILG kernel modifies structure
growth at wavenumbers $k > k_0^{-1} \sim 1$ Mpc$^{-1}$, potentially
producing a small shift in the observed BAO peak position at $z < 1$.
This is a falsifiable prediction for DESI DR5 and Euclid.

%%============================================================
\section{Conclusion}
\label{sec:conclusion}
%%============================================================

We have derived the BAO sound horizon from the RS framework.  The
characteristic scale $r_s \approx 147$ Mpc follows from the RS-derived
baryon-to-photon ratio, primordial spectrum, and Friedmann expansion.
The RS prediction agrees with BOSS measurements to $0.1\%$.  All
input quantities are derived from the $\varphi$-ladder with zero free
parameters; the sound horizon integral is identified as a hypothesis
with a clear falsifier for DESI/Euclid.

\begin{thebibliography}{99}

\bibitem{Eisenstein2005}
D.~J.~Eisenstein \textit{et al.}, ``Detection of baryon acoustic
oscillations in the correlation function of a sample of $70{,}000$
luminous red galaxies,'' \textit{Astrophys.\ J.}\ \textbf{633}, 560
(2005).

\bibitem{Alam2017}
S.~Alam \textit{et al.}\ (BOSS Collaboration), ``The clustering of
galaxies in the completed SDSS-III Baryon Oscillation Spectroscopic
Survey,'' \textit{Mon.\ Not.\ Roy.\ Astron.\ Soc.}\ \textbf{470},
2617 (2017).

\bibitem{Planck2018}
Planck Collaboration, ``Planck 2018 results. VI. Cosmological
parameters,'' \textit{Astron.\ Astrophys.}\ \textbf{641}, A6 (2020).

\bibitem{Eisenstein1998}
D.~J.~Eisenstein, W.~Hu, ``Baryonic features in the matter transfer
function,'' \textit{Astrophys.\ J.}\ \textbf{496}, 605 (1998).

\bibitem{Washburn2026a}
J.~Washburn, ``The Algebra of Reality: A Recognition Science Derivation
of Physical Law,'' \textit{Axioms} \textbf{15}(2), 90 (2026).

\bibitem{Washburn2026b}
J.~Washburn, ``Primordial Nucleosynthesis from Recognition Science,''
preprint (2026).

\end{thebibliography}

\end{document}
